\chapter{Sprint 8 : Tests, Validation et Déploiement}
\label{chap:sprint8-deploiement}

\section*{Introduction}
\addcontentsline{toc}{section}{Introduction}

Le dernier sprint d'un projet Agile est crucial : il marque la transition entre le développement et la mise en production. Ce sprint final est dédié à la validation complète de la solution, à la vérification de ses performances et de sa sécurité, ainsi qu'à son déploiement dans l'environnement de production de la BTL. Au-delà des aspects techniques, ce sprint inclut également la formation des utilisateurs et la documentation complète du système, garantissant ainsi une adoption réussie de la solution décisionnelle.

Dans ce chapitre, nous commencerons par définir les objectifs précis du sprint, puis nous présenterons la conception à travers des diagrammes UML. Nous détaillerons ensuite les tests de performance et de charge, la sécurité et l'audit, le déploiement des différentes composantes de la solution (backend, frontend, base de données), la formation des utilisateurs, et la documentation technique. Enfin, nous présenterons les livrables de ce sprint.

\section{Objectifs du Sprint}
\label{sec:sprint8-objectifs}

Le sprint 8 a pour objectifs principaux :

\begin{itemize}
    \item Réaliser des \textbf{tests de performance et de charge} pour valider la robustesse de la solution.
    \item Effectuer un \textbf{audit de sécurité} complet pour garantir la protection des données sensibles.
    \item \textbf{Déployer} la solution dans l'environnement de production de la BTL.
    \item \textbf{Former} les utilisateurs à l'utilisation de la solution décisionnelle.
    \item Produire une \textbf{documentation technique} complète pour la maintenance et l'évolution.
    \item Assurer un \textbf{transfert de compétences} vers les équipes IT de la BTL.
\end{itemize}

À la fin de ce sprint, la solution décisionnelle sera opérationnelle en production, les utilisateurs seront formés, et les équipes techniques disposeront de toute la documentation nécessaire pour assurer la maintenance et les évolutions futures.

\section{Conception du Sprint}
\label{sec:sprint8-conception}

La phase de conception de ce sprint est centrée sur l'architecture de déploiement, la configuration des environnements et les processus de validation.

\subsection{Diagramme de cas d'utilisation}
\label{subsec:sprint8-use-case}

Le diagramme de cas d'utilisation présente les fonctionnalités liées à l'administration système et au déploiement. Les acteurs identifiés sont :

\begin{itemize}
    \item \textbf{Administrateur système} : Responsable du déploiement, de la configuration et de la supervision de la solution.
    \item \textbf{Administrateur BTL} : Responsable de la gestion des utilisateurs et des droits d'accès en production.
    \item \textbf{Formateur} : Responsable de la formation des utilisateurs.
    \item \textbf{Utilisateur final} : Bénéficiaire de la solution déployée.
    \item \textbf{Système CI/CD} : Pipeline automatisé pour le déploiement continu.
\end{itemize}

\begin{figure}[H]
\centering
\includegraphics[scale=0.5]{images/diagramme_use_case_deploiement.png}
\caption{Diagramme de cas d'utilisation - Administration système et déploiement}
\label{fig:use-case-deploiement}
\end{figure}

\textbf{Description :} Le diagramme de cas d'utilisation présente les interactions entre les acteurs (Administrateur système, Administrateur BTL, Formateur, Utilisateur final, Système CI/CD) et les fonctionnalités de déploiement : configuration des environnements, tests de charge, audit de sécurité, déploiement de la solution, supervision, formation des utilisateurs, consultation de la documentation et gestion des utilisateurs.

Les cas d'utilisation identifiés pour ce sprint sont :

\begin{table}[H]
\centering
\begin{tabular}{|p{4cm}|p{3.5cm}|p{6cm}|}
\hline
\textbf{Cas d'utilisation} & \textbf{Acteurs} & \textbf{Description} \\
\hline
Configurer les environnements & Administrateur système & Paramétrer les environnements de développement, test et production. \\
\hline
Exécuter les tests de charge & Administrateur système & Simuler des charges pour valider la performance. \\
\hline
Auditer la sécurité & Administrateur système & Vérifier les vulnérabilités et la conformité. \\
\hline
Déployer la solution & Administrateur système & Installer et configurer la solution en production. \\
\hline
Superviser le système & Administrateur système & Surveiller les indicateurs de performance et les logs. \\
\hline
Former les utilisateurs & Formateur & Organiser des sessions de formation. \\
\hline
Consulter la documentation & Administrateur système & Accéder à la documentation technique. \\
\hline
Gérer les utilisateurs & Administrateur BTL & Gérer les comptes et les droits en production. \\
\hline
\end{tabular}
\caption{Cas d'utilisation du sprint 8}
\label{tab:use-cases-deploiement}
\end{table}

\subsection{Diagramme de classes}
\label{subsec:sprint8-classes}

Le diagramme de classes modélise la structure de la configuration du système et des environnements de déploiement.

Les principales classes identifiées sont :

\begin{itemize}
    \item \textbf{EnvironmentSettings} : Configuration d'un environnement spécifique.
        \begin{itemize}
            \item Attributs : environment\_name (dev, test, prod), database\_url, api\_url, frontend\_url, debug\_mode, log\_level.
            \item Méthodes : load\_settings(), validate\_settings(), get\_environment().
        \end{itemize}
    
    \item \textbf{DeploymentConfig} : Configuration du déploiement.
        \begin{itemize}
            \item Attributs : version, deploy\_timestamp, components (liste des composants), deployment\_strategy (blue-green, rolling, etc.).
            \item Méthodes : deploy(), rollback(), status().
        \end{itemize}
    
    \item \textbf{Component} : Représente un composant de la solution.
        \begin{itemize}
            \item Attributs : name, type (backend, frontend, database), version, status, health\_check.
            \item Méthodes : deploy\_component(), stop(), start(), restart().
        \end{itemize}
    
    \item \textbf{MonitoringConfig} : Configuration de la supervision.
        \begin{itemize}
            \item Attributs : metrics (liste des métriques), thresholds (seuils d'alerte), alert\_channels (canaux de notification).
            \item Méthodes : check\_health(), get\_metrics(), send\_alert().
        \end{itemize}
    
    \item \textbf{SecurityConfig} : Configuration de la sécurité.
        \begin{itemize}
            \item Attributs : ssl\_enabled, authentication\_type, password\_policy, session\_timeout, audit\_log\_enabled.
            \item Méthodes : validate\_security(), get\_security\_report().
        \end{itemize}
\end{itemize}

\begin{figure}[H]
\centering
\includegraphics[scale=0.5]{images/diagramme_classes_deploiement.png}
\caption{Diagramme de classes - Configuration et déploiement}
\label{fig:class-diagram-deploiement}
\end{figure}

\textbf{Description :} Le diagramme de classes présente la structure de la configuration du système avec les classes EnvironmentSettings (configuration d'environnement), DeploymentConfig (configuration du déploiement), Component (composant de la solution), MonitoringConfig (supervision) et SecurityConfig (sécurité). EnvironmentSettings est utilisé par DeploymentConfig, qui contient plusieurs Components, surveillés par MonitoringConfig, et SecurityConfig s'applique à l'ensemble.

Les relations entre les classes sont les suivantes :

\begin{itemize}
    \item \textbf{EnvironmentSettings} est utilisé par \textbf{DeploymentConfig}.
    \item \textbf{DeploymentConfig} contient plusieurs \textbf{Component}.
    \item \textbf{MonitoringConfig} surveille les \textbf{Component}.
    \item \textbf{SecurityConfig} est appliqué à l'ensemble des \textbf{Component}.
\end{itemize}

\subsection{Diagramme de séquence}
\label{subsec:sprint8-sequence}

Les diagrammes de séquence suivants illustrent les principaux flux d'exécution du déploiement.

\subsubsection{Processus de déploiement CI/CD}

\begin{figure}[H]
\centering
\includegraphics[scale=0.5]{images/diagramme_sequence_cicd.png}
\caption{Diagramme de séquence - Processus de déploiement CI/CD}
\label{fig:sequence-cicd}
\end{figure}

\textbf{Description :} Le diagramme de séquence du processus de déploiement CI/CD montre les interactions entre le Développeur, le dépôt Git, le pipeline CI/CD, l'environnement de test (Staging) et l'environnement de production. Le développeur pousse le code, le pipeline CI exécute les tests, build l'application, déploie en staging, effectue les tests d'acceptation, puis déploie en production avec stratégie (blue-green, canary, rolling) et effectue les health checks.

Les étapes du processus de déploiement CI/CD sont :

\begin{enumerate}
    \item Le développeur pousse son code vers le dépôt Git.
    \item Le pipeline CI (GitHub Actions, GitLab CI, Jenkins) est déclenché automatiquement.
    \item Les tests unitaires et d'intégration sont exécutés.
    \item Si les tests sont réussis, l'application est buildée (build).
    \item Un artefact (fichier image Docker, bundle) est créé.
    \item L'artefact est déployé dans l'environnement de test (staging).
    \item Les tests d'acceptation sont exécutés dans l'environnement de test.
    \item Si les tests d'acceptation sont réussis, le déploiement en production est approuvé.
    \item L'artefact est déployé en production avec la stratégie de déploiement définie (blue-green, canary, rolling).
    \item Les health checks sont effectués pour valider le bon fonctionnement.
    \item Si le déploiement échoue, un rollback automatique est déclenché.
    \item Un rapport de déploiement est généré et envoyé à l'administrateur.
\end{enumerate}

\subsubsection{Tests de charge}

\begin{figure}[H]
\centering
\includegraphics[scale=0.5]{images/diagramme_sequence_load_test.png}
\caption{Diagramme de séquence - Tests de charge}
\label{fig:sequence-load-test}
\end{figure}

\textbf{Description :} Le diagramme de séquence des tests de charge montre les interactions entre l'Administrateur, l'outil de test (K6/JMeter/Locust), l'API FastAPI, la base de données PostgreSQL et le frontend React. L'administrateur initialise l'outil, simule des utilisateurs simultanés, envoie des requêtes à l'API, mesure les temps de réponse, surveille la base de données et le frontend, collecte les résultats et génère un rapport de performance.

Les étapes des tests de charge sont :

\begin{enumerate}
    \item L'administrateur initialise l'outil de test de charge (ex: JMeter, Locust, K6).
    \item L'outil de test simule un nombre croissant d'utilisateurs simultanés.
    \item Des requêtes sont envoyées à l'API FastAPI avec différents scénarios.
    \item Les temps de réponse et les taux d'erreur sont mesurés.
    \item Les performances de la base de données PostgreSQL sont surveillées (CPU, I/O, connexions).
    \item Les performances du frontend React sont mesurées (temps de chargement, temps de rendu).
    \item Les résultats sont collectés et analysés.
    \item Un rapport de performance est généré avec les métriques et les recommandations.
\end{enumerate}

\section{Tests de performance et de charge}
\label{sec:sprint8-performance}

Les tests de performance et de charge visent à valider que la solution peut supporter un nombre d'utilisateurs et un volume de données réalistes.

\subsection{Objectifs des tests}

\begin{itemize}
    \item \textbf{Temps de réponse} : Vérifier que les requêtes API répondent en moins de 3 secondes en moyenne.
    \item \textbf{Capacité de charge} : Vérifier que la solution peut supporter 50 utilisateurs simultanés (concurrents).
    \item \textbf{Volume de données} : Vérifier le temps de chargement des rapports Power BI avec un volume de données équivalent à 3 ans d'historique.
    \item \textbf{Tolérance aux pics} : Vérifier le comportement de la solution sous une augmentation soudaine de la charge.
\end{itemize}

\subsection{Métriques de performance}

\begin{table}[H]
\centering
\begin{tabular}{|p{4cm}|p{4cm}|p{5cm}|}
\hline
\textbf{Métrique} & \textbf{Cible} & \textbf{Description} \\
\hline
Temps de réponse API & < 3 secondes & Temps moyen de réponse pour les appels API. \\
\hline
Temps de chargement dashboard & < 5 secondes & Temps d'affichage du dashboard React. \\
\hline
Temps de chargement Power BI & < 8 secondes & Temps d'affichage des rapports Power BI. \\
\hline
Taux d'erreur & < 1\% & Pourcentage de requêtes ayant échoué. \\
\hline
Nombre de requêtes par seconde & 100 RPS & Capacité maximale de l'API. \\
\hline
Temps de réponse ETL & < 30 minutes & Durée d'exécution des jobs ETL quotidiens. \\
\hline
\end{tabular}
\caption{Métriques de performance cibles}
\label{tab:performance-targets}
\end{table}

\subsection{Outils utilisés}

Les outils utilisés pour les tests de performance sont :

\begin{itemize}
    \item \textbf{K6} : Outil open source de test de charge pour les APIs.
    \item \textbf{Lighthouse} : Outil Google pour l'analyse de performance des pages web.
    \item \textbf{PostgreSQL Monitoring} : Outils de supervision de la base de données.
    \item \textbf{Power BI Performance Analyzer} : Outil intégré pour analyser la performance des rapports.
\end{itemize}

\subsection{Résultats des tests de performance}

\begin{table}[H]
\centering
\begin{tabular}{|p{4cm}|p{3.5cm}|p{3.5cm}|p{2.5cm}|}
\hline
\textbf{Scénario de test} & \textbf{Valeur mesurée} & \textbf{Cible} & \textbf{Statut} \\
\hline
API Dashboard - 10 utilisateurs & 0,8 seconde & < 3 secondes & ✓ \\
\hline
API Dashboard - 50 utilisateurs & 1,5 seconde & < 3 secondes & ✓ \\
\hline
API Dashboard - 100 utilisateurs & 2,8 secondes & < 5 secondes & ✓ \\
\hline
Chargement React Dashboard & 2,1 secondes & < 5 secondes & ✓ \\
\hline
Chargement Power BI Report & 4,5 secondes & < 8 secondes & ✓ \\
\hline
Job ETL quotidien & 18 minutes & < 30 minutes & ✓ \\
\hline
Taux d'erreur à 100 RPS & 0,3\% & < 1\% & ✓ \\
\hline
\end{tabular}
\caption{Résultats des tests de performance}
\label{tab:performance-results}
\end{table}

\section{Sécurité et audit}
\label{sec:sprint8-securite}

La sécurité est une priorité absolue pour une solution bancaire. Des mesures de sécurité et des mécanismes d'audit sont mis en place.

\subsection{Audit de sécurité}

L'audit de sécurité couvre les aspects suivants :

\begin{itemize}
    \item \textbf{Authentification et autorisation} : Vérification de la robustesse du système JWT, de la gestion des rôles et des permissions.
    \item \textbf{Chiffrement} : Vérification du chiffrement des données en transit (HTTPS) et au repos (base de données).
    \item \textbf{Protection contre les injections} : Vérification que toutes les requêtes SQL sont paramétrées.
    \item \textbf{Protection XSS} : Vérification que les entrées utilisateur sont correctement échappées.
    \item \textbf{Protection CSRF} : Vérification que les tokens CSRF sont utilisés.
    \item \textbf{Gestion des secrets} : Vérification que les mots de passe, clés API et tokens sont stockés de manière sécurisée (variables d'environnement, Vault).
    \item \textbf{Conformité RGPD} : Vérification que les données personnelles sont traitées conformément à la réglementation.
\end{itemize}

\subsection{Mécanismes d'audit}

Les mécanismes d'audit incluent :

\begin{itemize}
    \item \textbf{Logs d'accès} : Enregistrement de toutes les connexions utilisateur (date, heure, IP, identifiant).
    \item \textbf{Logs d'action} : Enregistrement des actions sensibles (création/modification/suppression d'utilisateur, consultation de données sensibles).
    \item \textbf{Logs d'erreur} : Enregistrement des erreurs système et des tentatives d'accès non autorisées.
    \item \textbf{Conservation des logs} : Les logs sont conservés pendant une durée définie (ex: 6 mois) pour permettre les audits rétrospectifs.
\end{itemize}

\subsection{Plan de reprise d'activité (PRA)}

Un plan de reprise d'activité est défini :

\begin{itemize}
    \item \textbf{Sauvegardes régulières} : Sauvegarde quotidienne de la base de données PostgreSQL.
    \item \textbf{Réplication} : Réplication de la base de données vers un site secondaire.
    \item \textbf{Restauration} : Procédure de restauration en cas de panne (RTO < 4 heures, RPO < 1 heure).
    \item \textbf{Redondance} : Déploiement des composants avec redondance (plusieurs instances).
\end{itemize}

\begin{figure}[H]
\centering
\includegraphics[scale=0.5]{images/securite_audit_pra.png}
\caption{Mécanismes de sécurité et plan de reprise d'activité}
\label{fig:securite-audit}
\end{figure}

\textbf{Description :} La sécurité et l'audit reposent sur plusieurs piliers : l'audit de sécurité (authentification JWT, chiffrement HTTPS, protection SQL/XSS/CSRF, gestion des secrets, conformité RGPD), les mécanismes d'audit (logs d'accès, logs d'action, logs d'erreur, conservation des logs), et le plan de reprise d'activité (sauvegardes quotidiennes, réplication, restauration RTO<4h/RPO<1h, redondance).

\section{Déploiement}
\label{sec:sprint8-deploiement}

Le déploiement de la solution dans l'environnement de production de la BTL est l'étape finale du projet.

\subsection{Architecture de déploiement}

L'architecture de déploiement est composée de plusieurs couches :

\begin{itemize}
    \item \textbf{Couche infrastructure} : Serveurs physiques ou virtuels (VM) hébergés dans le datacenter de la BTL ou dans le cloud.
    \item \textbf{Couche applicative} : Conteneurs Docker orchestrés par Kubernetes (ou Docker Compose pour une infrastructure plus légère).
    \item \textbf{Couche réseau} : Reverse proxy (Nginx) pour le routage et le load balancing.
    \item \textbf{Couche données} : PostgreSQL avec réplication et sauvegardes automatisées.
\end{itemize}

\begin{figure}[H]
\centering
\includegraphics[scale=0.5]{images/architecture_deploiement.png}
\caption{Architecture de déploiement de la solution}
\label{fig:deployment-architecture}
\end{figure}

\textbf{Description :} L'architecture de déploiement est organisée en quatre couches : Infrastructure (serveurs physiques/virtuels dans le datacenter BTL ou cloud), Applicative (conteneurs Docker orchestrés par Kubernetes ou Docker Compose), Réseau (reverse proxy Nginx pour routage et load balancing), et Données (PostgreSQL avec réplication et sauvegardes automatisées).

\subsection{Déploiement du Backend FastAPI}
\label{subsec:deploiement-backend}

Le déploiement du backend FastAPI comprend :

\begin{itemize}
    \item \textbf{Conteneurisation} : L'application FastAPI est conteneurisée avec Docker.
    \item \textbf{Base de données} : Les connexions à PostgreSQL sont configurées via des variables d'environnement.
    \item \textbf{Passerelle API} : Un reverse proxy (Nginx) est configuré pour router les requêtes vers l'API FastAPI.
    \item \textbf{Load Balancing} : Plusieurs instances de l'API peuvent être déployées pour la montée en charge.
    \item \textbf{Supervision} : Les endpoints de health check sont configurés pour la supervision.
\end{itemize}

\subsection{Déploiement du Frontend React}
\label{subsec:deploiement-frontend}

Le déploiement du frontend React comprend :

\begin{itemize}
    \item \textbf{Build} : L'application React est buildée en version de production (npm run build).
    \item \textbf{Servir les fichiers statiques} : Les fichiers HTML, CSS, JS sont servis par un serveur web (Nginx).
    \item \textbf{Configuration Nginx} : Nginx est configuré pour servir les fichiers statiques, gérer le cache, et rediriger les requêtes API vers le backend.
    \item \textbf{CDN} : Les fichiers peuvent être déployés sur un CDN pour améliorer les performances.
\end{itemize}

\subsection{Déploiement de PostgreSQL}
\label{subsec:deploiement-database}

Le déploiement de PostgreSQL comprend :

\begin{itemize}
    \item \textbf{Installation} : PostgreSQL est installé sur le serveur de base de données.
    \item \textbf{Configuration} : PostgreSQL est configuré pour les performances (shared\_buffers, work\_mem, etc.).
    \item \textbf{Sécurité} : L'authentification est configurée (md5 ou SCRAM), les connexions SSL sont activées.
    \item \textbf{Sauvegardes} : Des sauvegardes automatiques sont planifiées (pg\_dump, WAL archiving).
    \item \textbf{Réplication} : La réplication est configurée pour la haute disponibilité (si nécessaire).
\end{itemize}

\subsection{Déploiement avec Kubernetes (optionnel)}

Si l'infrastructure le permet, un déploiement avec Kubernetes offre :

\begin{itemize}
    \item \textbf{Orchestration} : Gestion automatique des conteneurs (démarrage, arrêt, redémarrage).
    \item \textbf{Mise à l'échelle automatique} : Ajustement automatique du nombre de réplicas en fonction de la charge.
    \item \textbf{Rolling updates} : Mises à jour sans interruption de service.
    \item \textbf{Service discovery} : Découverte automatique des services.
\end{itemize}

\section{Formation des utilisateurs}
\label{sec:sprint8-formation}

La formation des utilisateurs est essentielle pour garantir l'adoption de la solution décisionnelle.

\subsection{Guides d'utilisation}

Des guides d'utilisation sont produits pour chaque type d'utilisateur :

\begin{itemize}
    \item \textbf{Guide du Directeur} :
        \begin{itemize}
            \item Comment consulter les tableaux de bord Power BI.
            \item Comment comparer les agences.
            \item Comment utiliser le chatbot pour des analyses stratégiques.
            \item Comment visualiser les prédictions.
        \end{itemize}
    \item \textbf{Guide du Responsable d'agence} :
        \begin{itemize}
            \item Comment suivre les KPI de son agence.
            \item Comment utiliser le chatbot pour des questions opérationnelles.
            \item Comment visualiser les prédictions pour son agence.
            \item Comment exporter les rapports.
        \end{itemize}
    \item \textbf{Guide de l'Administrateur} :
        \begin{itemize}
            \item Comment gérer les utilisateurs et les droits d'accès.
            \item Comment superviser les jobs ETL.
            \item Comment configurer les alertes.
            \item Comment assurer la maintenance du système.
        \end{itemize}
\end{itemize}

\subsection{Sessions de formation}

Des sessions de formation sont organisées :

\begin{table}[H]
\centering
\begin{tabular}{|p{3cm}|p{4.5cm}|p{6cm}|}
\hline
\textbf{Session} & \textbf{Public cible} & \textbf{Contenu} \\
\hline
Session 1 & Directeurs & Présentation de la solution, démonstration des dashboards Power BI, utilisation du chatbot et des prédictions. \\
\hline
Session 2 & Responsables d'agence & Présentation de la solution, démonstration des dashboards détaillés, utilisation du chatbot, visualisation des prédictions. \\
\hline
Session 3 & Administrateurs & Formation technique : gestion des utilisateurs, supervision ETL, maintenance, résolution des problèmes. \\
\hline
Session 4 & Équipe IT & Formation technique approfondie : architecture, déploiement, API, intégration Power BI. \\
\hline
\end{tabular}
\caption{Sessions de formation}
\label{tab:formation-sessions}
\end{table}

\subsection{Matériel de formation}

Le matériel de formation comprend :

\begin{itemize}
    \item \textbf{Présentations} : Slides pour chaque session de formation.
    \item \textbf{Vidéos tutorielles} : Vidéos courtes démontrant les fonctionnalités clés.
    \item \textbf{Fiches pratiques} : Fiches synthétiques pour les actions fréquentes.
    \item \textbf{FAQ} : Foire aux questions avec les réponses aux questions courantes.
\end{itemize}

\subsection{Suivi de la formation}

Un suivi de la formation est effectué :

\begin{itemize}
    \item \textbf{Questionnaires de satisfaction} : Recueil des retours après chaque session.
    \item \textbf{Évaluation des compétences} : Test des compétences acquises.
    \item \textbf{Assistance post-formation} : Période de support pour les questions après la formation.
\end{itemize}

\section{Documentation technique}
\label{sec:sprint8-documentation}

La documentation technique est essentielle pour la maintenance et l'évolution de la solution.

\subsection{Documentation d'architecture}

La documentation d'architecture comprend :

\begin{itemize}
    \item \textbf{Architecture globale} : Schéma de l'architecture de la solution.
    \item \textbf{Diagrammes UML} : Cas d'utilisation, classes, séquence.
    \item \textbf{Flux de données} : Description des flux entre les composants.
    \item \textbf{Modèle de données} : Description du Data Warehouse (tables, relations).
\end{itemize}

\subsection{Documentation API}

La documentation API comprend :

\begin{itemize}
    \item \textbf{Swagger/OpenAPI} : Documentation interactive des endpoints.
    \item \textbf{Guides d'utilisation} : Comment utiliser chaque endpoint.
    \item \textbf{Exemples} : Exemples de requêtes et de réponses.
    \item \textbf{Authentification} : Comment s'authentifier pour utiliser l'API.
\end{itemize}

\subsection{Documentation de déploiement}

La documentation de déploiement comprend :

\begin{itemize}
    \item \textbf{Prérequis} : Configuration matérielle et logicielle requise.
    \item \textbf{Procédure d'installation} : Étapes d'installation et de configuration.
    \item \textbf{Configuration des environnements} : Paramètres pour les environnements de développement, test et production.
    \item \textbf{Procédure de mise à jour} : Comment mettre à jour la solution.
    \item \textbf{Procédure de sauvegarde et restauration} : Comment sauvegarder et restaurer les données.
\end{itemize}

\subsection{Documentation de maintenance}

La documentation de maintenance comprend :

\begin{itemize}
    \item \textbf{Procédures de surveillance} : Comment surveiller le système.
    \item \textbf{Gestion des erreurs} : Comment identifier et résoudre les erreurs courantes.
    \item \textbf{Procédures de scaling} : Comment faire évoluer l'infrastructure.
    \item \textbf{Plan de reprise d'activité} : Procédure en cas de sinistre.
\end{itemize}

\section{Livrables du Sprint}
\label{sec:sprint8-livrables}

À l'issue de ce sprint, les livrables suivants sont produits :

\begin{table}[H]
\centering
\begin{tabular}{|p{4cm}|p{9cm}|}
\hline
\textbf{Livrable} & \textbf{Description} \\
\hline
Rapport de tests de performance & Résultats détaillés des tests de charge et de performance. \\
\hline
Rapport d'audit de sécurité & Résultats de l'audit de sécurité et recommandations. \\
\hline
Solution déployée en production & Backend FastAPI, Frontend React, Base de données PostgreSQL. \\
\hline
Documentation de déploiement & Procédure d'installation et de configuration. \\
\hline
Guides utilisateurs & Guides pour Directeur, Responsable, Administrateur. \\
\hline
Matériel de formation & Présentations, vidéos, fiches pratiques, FAQ. \\
\hline
Documentation technique & Architecture, API, Maintenance. \\
\hline
Rapport de formation & Résumé des sessions de formation et retours utilisateurs. \\
\hline
\end{tabular}
\caption{Livrables du sprint 8}
\label{tab:sprint8-livrables}
\end{table}



\textbf{Description :} Les livrables du sprint 8 comprennent : le rapport de tests de performance (résultats détaillés des tests de charge), le rapport d'audit de sécurité (recommandations), la solution déployée en production (FastAPI, React, PostgreSQL), la documentation de déploiement (procédure d'installation), les guides utilisateurs (Directeur, Responsable, Administrateur), le matériel de formation (présentations, vidéos, fiches pratiques, FAQ), la documentation technique (architecture, API, maintenance) et le rapport de formation (résumé des sessions et retours utilisateurs).

\section*{Conclusion}
\addcontentsline{toc}{section}{Conclusion}

Ce huitième et dernier sprint a marqué l'aboutissement de notre projet de solution décisionnelle pour la BTL. Les tests de performance et de charge ont validé la robustesse de la solution, avec des temps de réponse conformes aux exigences et une capacité à supporter la charge prévue. L'audit de sécurité a confirmé la protection des données sensibles et la conformité aux bonnes pratiques.

Le déploiement en production a été réalisé avec succès, avec une architecture conteneurisée garantissant la portabilité et la maintenabilité de la solution. La formation des utilisateurs, organisée en plusieurs sessions, a permis de familiariser les responsables de la BTL avec les fonctionnalités de la solution décisionnelle : tableaux de bord Power BI, chatbot intelligent, et module de prédiction.

La documentation technique complète produite assure la pérennité du projet, facilitant la maintenance et les évolutions futures par les équipes IT de la BTL.

Ce sprint final clôture le projet dans son ensemble. La solution décisionnelle intelligente est désormais opérationnelle, offrant à la BTL un outil puissant pour suivre, évaluer et optimiser les performances de ses agences.